%% Chapter 9, Section 9.4 — Discussion and Bibliography
%% A First Course in Monte Carlo Methods
%% Sanz-Alonso & Al-Ghattas

\section{Discussion and Bibliography}
\label{sec:9.4}

Sequential Monte Carlo methods are overviewed from an algorithmic viewpoint in
\cite{58,59} and from a mathematical perspective in \cite{54,49}. The term ``bootstrap''
particle filter was coined in the seminal work \cite{94} in reference to the bootstrap
\cite{70}, a standard statistical procedure that also involves sampling with replacement.
The optimal particle filter is discussed, and further references given, in \cite{59}; see
section~IID. While the term ``optimal'' particle filter is common in the literature, the
sense in which this filter is optimal is indeed rather weak. We refer to \cite{49}, Chapter~10,
which utilizes instead the term ``guided'' particle filter, for further discussion on
this topic; in particular, \cite{49}, Theorem~10.1, formalizes the criterion of
\emph{local optimality} and the ensuing discussion provides compelling criticisms of this
notion of optimality.

The proof of convergence of the bootstrap particle filter presented here originates in
\cite{190} and can also be found in \cite{208}. We refer to \cite{208}, Chapter~12, for a
similar proof of convergence for a Gaussianized optimal particle filter, and to \cite{118}
for a convergence analysis of the optimal particle filter. Throughout much of this
chapter, we considered for simplicity the case of Gaussian additive noise and linear
observation operator. More general convergence and stability results for sequential Monte
Carlo methods can be found, for instance, in \cite{54} and \cite{49}, Chapter~11.

For problems in which the dynamics model is nonlinear and the state space of the hidden
process is low dimensional, particle filters are enormously successful. Generalizing them
so that they work for the high-dimensional problems that arise, for example, in
geophysical applications, provides a major challenge, since in those settings the particle
weights typically concentrate on one, or a small number, of particles --- see for instance
\cite{24, 216, 215, 3}. Weight collapse is intimately related to the fact that, as
discussed in Chapter~3, importance sampling weights have large variance when the target
and the proposal distributions are far apart. As noted in \cite{217, 3}, an advantage of
the optimal particle filter over the bootstrap particle filter is that the optimal filter
utilizes a proposal that is closer to the target, which helps alleviate weight collapse
over each iteration of the filter. Resampling techniques to turn weighted samples into
unweighted ones are essential to the success of particle filters over multiple iterations.
We refer to \cite{49}, Chapter~9, for an in-depth exploration and comparison of resampling
schemes in sequential Monte Carlo.

An important algorithmic idea in sequential Monte Carlo, as in many other algorithms
studied in these notes, is to introduce auxiliary variables to accelerate computation
\cite{186, 118}. In addition, exploiting decay of correlations through \emph{localization}
is often needed when implementing particle filters for online estimation of
discretizations of spatial fields. A review of local particle filters can be found in
\cite{73}. The paper \cite{190} investigates, from a theoretical viewpoint, whether
localization can help to beat the curse of dimension. Attempts to bridge particle filters
with ensemble Kalman filters to alleviate the curse of dimension include \cite{76, 218},
and the relation between the collapse of ensemble and particle methods is investigated in
the paper \cite{165}, which also emphasizes the importance of localization. Further
theoretical investigations of localization for ensemble Kalman methods include \cite{6,
5}. We close by pointing out that sequential Monte Carlo methods can be combined in
several ways with algorithms studied in previous chapters, and we refer to \cite{10} for a
representative, important example.
